\documentclass[11pt,a4paper]{article}
\usepackage[margin=1in]{geometry}
\usepackage{amsmath,amssymb}
\usepackage{graphicx}
\usepackage{float}
\usepackage{enumitem}
\usepackage{microtype}
\usepackage{caption}
\usepackage{tabularx}
\usepackage{hyperref}

\captionsetup{font=small,labelfont=bf,skip=6pt}

\setlist[itemize]{leftmargin=1.2cm}
\setlist[enumerate]{leftmargin=1.2cm}

\newcommand{\problembox}[3]{%
  \par\addvspace{0.42em}%
  \noindent
  \setlength{\fboxsep}{7.5pt}%
  \setlength{\fboxrule}{0.95pt}%
  \fbox{%
    \begin{minipage}{\dimexpr\linewidth-2\fboxsep-2\fboxrule\relax}
    \begin{tabularx}{\linewidth}{@{}>{\raggedright\arraybackslash\bfseries}p{2.1em} @{\hspace{0.9em}} >{\raggedright\arraybackslash}X @{\hspace{0.9em}} >{\raggedleft\arraybackslash\bfseries}p{2.4em}@{}}
    #1 & #2 & #3
    \end{tabularx}
    \end{minipage}%
  }%
  \par\addvspace{0.32em}%
}

\title{\textbf{Flight to the Moon}}
\author{\normalsize Y26 (2026) --- English translation}
\date{}

\begin{document}
\maketitle

In this problem, we will consider orbits that allow a flight to the Moon and a return back. In the simplest cases, such a flight is typically modeled as an elliptical orbit with a major axis equal to $r_m$, the distance from the Earth to the Moon. However, such an orbit is not optimal or safe. The main disadvantage is that the spacecraft must be captured by the Moon's gravity. During this capture, engines must be used to enter orbit; a failure at this stage could result in a crash on the lunar surface.

This safety issue is addressed by a \textbf{free-return trajectory}. The Moon acts as a gravitational slingshot, creating a "figure-eight" trajectory. This was used for the Soviet \textit{Luna-3} probe and for the emergency return of the \textit{Apollo 13} mission.

\begin{figure}[H]
    \centering
    \includegraphics[width=0.9\textwidth]{T4_EN_assets/fig1.jpg}
    \caption{Chronology of the Apollo 13 mission.}
\end{figure}

The following physical constants may be required:
\begin{itemize}
    \item Gravitational constant: $G = 6.67 \cdot 10^{-11}~\frac{\text{m}^3}{\text{kg} \cdot \text{s}^2}$
    \item Mass of the Earth: $M_e = 6 \cdot 10^{24}~\text{kg}$
    \item Mass of the Moon: $M_m = M_e / 81$
    \item Radius of the Earth: $R_e = 6400~\text{km}$
    \item Radius of the Moon: $R_m = 1700~\text{km}$
    \item Earth-Moon distance: $r_m = 3.85 \cdot 10^8~\text{m}$
    \item Orbital period of the Moon: $T_m = 27.4~\text{days}$
    \item Orbital speed of the Moon: $v_m = 2\pi r_m / T_m$
\end{itemize}

\section*{Part A. Takeoff from Earth (3.0 points)}

Most of a rocket's mass is first-stage fuel required to reach orbital velocity and overcome the atmosphere. We will consider a simplified vertical takeoff along the $y$-axis. The rocket has initial mass $m_0$ and is subject to gravity and drag $\vec{F} = -kv\vec{v}$. It expels gas at a constant mass flow rate $\mu$ with relative speed $u$ relative to the rocket.

Numerical values: $g = 9.8~\text{m/s}^2$, $\mu = 2700~\text{kg/s}$, $u = 4.5~\text{km/s}$, $k = 0.25~\text{kg/m}$ (near surface), $a = 6600~\text{km}$ (orbital radius), $m_0 = 500~\text{T}$.

\problembox{A1}{Write the expression for $\ddot{y}$ in terms of $\dot{y}$, $k$, $m_0$, $\mu$, $u$, $g$, and $t$.}{0.50}

\problembox{A2}{Neglecting drag, find the speed $v(m)$ and coordinate $y(m)$ at the moment the mass is $m$. Express in terms of $m$, $g$, $u$, $\mu$, and $m_0$.}{1.00}

\problembox{A3}{Assuming $k$ is constant with height, find numerically the height $H_c$ where the drag force equals the force of gravity.}{0.80}

\problembox{A4}{Estimate the ratio $\beta_1$ of initial to final mass required to reach circular orbital speed $v_a$ at radius $a$. Use $\beta_1 = 15$ for later parts.}{0.70}

\section*{Part B. Free Return Trajectory (6.0 points)}

\problembox{B1}{Find the distance $r_c$ from the Moon where the gravitational forces of the Earth and Moon are equal in magnitude.}{0.30}

The sphere of radius $r_c$ is the "Lunar zone of attraction". Inside, we only consider Lunar gravity ($\Delta U = \Delta U_m$); outside, only Earth gravity ($\Delta U = \Delta U_e$). The ship starts from a circular orbit of radius $a$, accelerating to $v_0$. Let $b$ be the minimum distance to the Moon's center.

\begin{figure}[H]
    \centering
    \includegraphics[width=0.7\textwidth]{T4_EN_assets/fig2.jpg}
    \caption{Trajectory geometry.}
\end{figure}

\problembox{B2}{In the Moon's frame, write the specific angular momentum $L_0^m$ and speed $v_2$ at distance $r_c$. Use $s = \sin(\alpha + \theta)$, $s' = \sin \alpha$, and $c = \cos \theta$.}{0.50}

\problembox{B3}{Using energy conservation in the Moon's frame, find $v_1/v_m$ in terms of $s$, $s'$, $c$, $b/r_c$, and $v_{\text{II}}^m/v_m$.}{1.20}

\problembox{B4}{Expand $v_1$ to first order in $b/r_c$ to show $v_1 \approx \frac{v_m c}{s} + B_0 \frac{b}{r_c}$. Express $B_0$.}{1.00}

\problembox{B5}{Using $\alpha \ll 1$, show that $v_1 \approx \frac{v_m}{\tan \theta} + A\alpha + B\frac{b}{r_c}$. Express $A$ and $B$.}{0.60}

\problembox{B6}{In the Earth's frame, find $\alpha$ to first order in $r_c/r_m$ and $a/r_m$ in terms of $a, r_c, r_m, v_1, v_0, \theta$.}{0.60}

\problembox{B7}{Neglecting terms of order $a/r_m$ and $r_c/r_m$, express $v_0$ in terms of $v_1$ and $v_{\text{II}}^e$.}{0.40}

\problembox{B8}{For a symmetric orbit, use the conic section equation to find an equation for $\theta$ using $v_1, \alpha, v_m, G, M_m, b, r_c$.}{1.00}

\problembox{B9}{Find $\beta_2$, the mass ratio for acceleration from $v_a$ to $v_0$. Calculate its value using data from Part A.}{0.40}

\section*{Part C. Lunar Landing (1.0 points)}

Assume a close approach with $b = 1800~\text{km}$. Since $h = b - R_m \ll R_m$, fuel is only used to brake to zero speed at the pericenter.

\problembox{C1}{With $v_2 = 1.5~\text{km/s}$, find the mass reduction factor $\beta_3$ during landing.}{0.50}

\problembox{C2}{Find the initial mass $m$ required to deliver a $40~\text{ton}$ cargo to the Moon.}{0.50}

\end{document}